\documentclass[11pt]{article}
\usepackage[margin=1in]{geometry}
\usepackage{amsmath,amssymb}
\usepackage{xcolor}
\usepackage{listings}
\usepackage{hyperref}
\hypersetup{colorlinks=true,linkcolor=blue,urlcolor=blue}
\usepackage{enumitem}
\usepackage{booktabs}

\definecolor{codebg}{rgb}{0.96,0.96,0.94}
\definecolor{pathcol}{rgb}{0.55,0.0,0.0}
\newcommand{\pp}[1]{\textcolor{pathcol}{\texttt{#1}}}

\lstset{
  basicstyle=\ttfamily\footnotesize,
  backgroundcolor=\color{codebg},
  breaklines=true,
  frame=single,
  language=C++,
  showstringspaces=false,
  columns=fullflexible,
}

\title{\textbf{Resolution, Smearing, and Unfolding in the Run-24 $p{+}p$\\
EMCal $\pi^0/\eta$ Non-Linearity Study}\\[4pt]
\large A Reading and Teaching Report on B.~Seidlitz's Working Directories}
\author{Prepared for: P.~Lewis \quad (read-only inspection of collaborator areas)}
\date{2026-06-15}

\begin{document}
\maketitle

\begin{abstract}
This report explains the three concepts that drive the sPHENIX Run-24 $p{+}p$
$\pi^0/\eta$ EMCal calibration, non-linearity, and PPG12 isolated-photon work:
\textbf{resolution}, \textbf{smearing}, and \textbf{unfolding}. Each concept is
first taught from first principles, then tied line-by-line to the actual code
that implements it in Blair Seidlitz's working areas
(\pp{/sphenix/u/bseidlitz/work/} and its GPFS mirror
\pp{/sphenix/user/bseidlitz/} --- the two are the same content). Every file and
directory inspected is named explicitly so the report can be retraced. A
companion plain-text inspection log, \texttt{RESO\_SMEAR\_UNFOLD\_LOG.txt},
lists the exact paths, sizes, and what was read.
\end{abstract}

\tableofcontents

%==============================================================
\section{Background and context of the study}
%==============================================================

The physics goal is a clean measurement of $\pi^0$ and $\eta$ mesons (and
isolated photons, PPG12) in $\sqrt{s}=200$~GeV $p{+}p$ collisions recorded by
sPHENIX in Run~24. Both mesons are reconstructed through their two-photon decay
channel,
\[
  \pi^0 \to \gamma\gamma \qquad (m_{\pi^0}=0.134977~\text{GeV}),\qquad
  \eta  \to \gamma\gamma \qquad (m_{\eta}=0.547862~\text{GeV}),
\]
by pairing electromagnetic clusters in the EMCal (CEMC) and forming the
two-photon invariant mass
\[
  M_{\gamma\gamma}=\sqrt{2E_1E_2\,(1-\cos\theta_{12})}.
\]
The measured peak position $M_{\rm fit}$ probes the \emph{energy scale} and its
\emph{non-linearity} (does $M_{\rm fit}/m_{\rm PDG}$ stay flat versus $p_T$?),
while the measured peak width $\sigma(M)$ probes the \emph{energy resolution}.

To compare data with simulation fairly, the Geant4 MC must reproduce the
detector's true performance. In practice the MC is slightly \emph{too good}
(its clusters are sharper than the real detector). The study therefore:
\begin{enumerate}[nosep]
  \item measures the resolution and energy scale in data via the meson peaks,
  \item \textbf{smears} the MC so that simulated peak widths and positions match
        data,
  \item propagates that tuned MC into efficiency and \textbf{unfolding}
        corrections for the final $p_T$ spectra / cross-sections (PPG12).
\end{enumerate}

The directories examined for this report:
\begin{itemize}[nosep]
  \item \pp{/sphenix/u/bseidlitz/work/resolutionStudy/} --- standalone toy-MC
        resolution/smearing demonstration.
  \item \pp{/sphenix/u/bseidlitz/work/emcPerf/calo\_emc\_pi0\_tbt/} --- the
        production tower-by-tower $\pi^0/\eta$ calibration module
        (\texttt{pi0EtaByEta.cc}), which contains the real smearing routine.
  \item \pp{/sphenix/u/bseidlitz/work/emcPerf/macros/figMaker/} ---
        \texttt{smear\_compare.C} and \texttt{function\_compare\_wide.root},
        where the data/MC smearing functions are derived and stored.
  \item \pp{/sphenix/u/bseidlitz/work/adjMCCalib/} --- MC energy-scale
        fine-tuning CDB trees (\texttt{convSwitch.C}).
  \item \pp{/sphenix/u/bseidlitz/work/ppg12/toymcunfold/macro/} ---
        \texttt{resolutionsim.C} and \texttt{jet\_pt\_unfolding.C}, the
        RooUnfold examples.
\end{itemize}

%==============================================================
\section{Resolution --- what it is and how it is defined here}
%==============================================================

\subsection{Teaching: the meaning of resolution}
\textbf{Energy resolution} is the fractional uncertainty with which the
calorimeter measures a particle's energy. If a fixed-energy photon is sent into
the EMCal many times, the reconstructed energies form a roughly Gaussian
distribution of width $\sigma_E$ about a mean $E$. The resolution is the ratio
\[
  \boxed{\;\frac{\sigma_E}{E}=\frac{a}{\sqrt{E}}\;\oplus\;c\;
          =\sqrt{\Big(\tfrac{a}{\sqrt{E}}\Big)^{2}+c^{2}}\;}
\]
where $\oplus$ denotes addition in quadrature. The two physically distinct
terms are:
\begin{description}[nosep]
  \item[Stochastic term $a/\sqrt{E}$:] from shower-sampling and
        photo-statistics. It \emph{shrinks} as energy rises, so high-energy
        showers are measured proportionally better.
  \item[Constant term $c$:] from calibration spread, light-collection
        non-uniformity, leakage, and dead material. It does \emph{not} shrink
        with energy and therefore dominates at high $E$.
\end{description}
Because $\pi^0/\eta$ widths come from \emph{two} photons, the meson relative
width $\sigma(M)/M$ inherits this same $a/\sqrt{E}\oplus c$ shape versus $p_T$
and is the experimental handle used throughout these directories.

\subsection{Where it is defined in code}
The cleanest written-out definition is the standalone toy in\\
\pp{/sphenix/u/bseidlitz/work/resolutionStudy/eta\_decay\_simulation\_with\_smearing.C}.
There the constant and stochastic terms are literal constants and are combined
in quadrature exactly as above:

\begin{lstlisting}
// energyResolution = constant term c = 12%
const double energyResolution = 0.12;
...
// stochastic term a/sqrt(E) with a = 15%
double energyDepSmear1 = 0.15 / TMath::Sqrt(photon1.E());
// quadrature sum  c (+) a/sqrt(E)  -> width of the Gaussian smear
double smearFactor1 = 1 + randGen.Gaus(0,
      TMath::Sqrt(energyResolution*energyResolution
                  + energyDepSmear1*energyDepSmear1));
\end{lstlisting}

This is exactly the curve labelled in the plotting macro
\pp{.../resolutionStudy/plot.C} as

\begin{lstlisting}
myText(0.65,0.30+0.5,1,"Toy MC   15/#sqrt{E} #oplus 12%");
\end{lstlisting}

i.e. $\sigma_E/E = 15\%/\sqrt{E}\,\oplus\,12\%$. The companion file
\pp{.../resolutionStudy/etaDecaySim.C} is an earlier variant with
$a=20\%$, $c=15\%$. The output \texttt{output.root} /
\texttt{meson\_graphs.root} feed \texttt{plot.C}, which fits each $p_T$ slice
with a Gaussian, forms $\sigma^{\eta}/M^{\eta}$, and overlays the toy curve on
the $p{+}p$ data points (\texttt{width\_vs\_pt.png}). \textbf{This is the
operational definition of ``resolution'' in the study: the $p_T$ dependence of
the fitted meson relative width.}

%==============================================================
\section{Smearing --- what it is and how it is done here}
%==============================================================

\subsection{Teaching: the meaning of smearing}
\textbf{Smearing} is the deliberate degradation of a simulated quantity so that
the MC reproduces the (worse) resolution of real data. The truth-level MC
photon energy $E$ is replaced by
\[
  E \;\longrightarrow\; E\cdot f,\qquad
  f \sim \mathrm{Gaus}\!\big(1,\;\sigma_E/E\big),
\]
a random multiplicative factor drawn from a Gaussian centred at 1 with width
equal to the desired \emph{extra} resolution. Smearing is the inverse mindset
to resolution: resolution \emph{measures} the blur; smearing \emph{injects} the
blur into an otherwise-too-perfect simulation. Position (angle) smearing works
the same way but is \emph{additive} on $\eta$ and $\phi$ and centred at zero.

\subsection{Toy-level smearing}
In the toy
\pp{.../resolutionStudy/eta\_decay\_simulation\_with\_smearing.C} the photon
4-vector is scaled by the Gaussian factor while keeping its direction fixed
(lines~69--70):
\begin{lstlisting}
photon1.SetPxPyPzE(photon1.Px()*smearFactor1, photon1.Py()*smearFactor1,
                   photon1.Pz()*smearFactor1, photon1.E()*smearFactor1);
\end{lstlisting}

\subsection{Production smearing: \texttt{pi0EtaByEta::smearPhot}}
The real, analysis-grade smearing lives in\\
\pp{/sphenix/u/bseidlitz/work/emcPerf/calo\_emc\_pi0\_tbt/pi0EtaByEta.cc},
function \texttt{smearPhot()} (lines~863--927). Key points:
\begin{itemize}[nosep]
  \item The energy- and position-resolution widths are \emph{not} hard-coded
        constants; they are read as ROOT \texttt{TF1} functions of energy from\\
        \pp{.../emcPerf/macros/figMaker/output/function\_compare\_wide.root}
        (opened in \texttt{Init()} at line~206, variable \texttt{kSmearParFile}).
  \item Energy smear: \texttt{Efac = rnd->Gaus(1.0, E\_res)} with
        \texttt{E\_res = f\_energy\_smear[idx]->Eval(E)}; clamped non-negative.
  \item Position smear: \texttt{Pfac\_eta/phi = rnd->Gaus(0.0, P\_res)} added to
        $\eta,\phi$ with $\phi$ wrapped to $(-\pi,\pi]$.
  \item Defaults if the file is missing: \texttt{E\_res = 0.08} (8\%),
        \texttt{P\_res = 0.0001}.
\end{itemize}

\begin{lstlisting}
float Efac     = rnd->Gaus(1.0f, E_res);     // energy smear factor
float Pfac_phi = rnd->Gaus(0.0f, P_res);     // angular smear (phi)
float Pfac_eta = rnd->Gaus(0.0f, P_res);     // angular smear (eta)
...
const float E_s = E * Efac;
sp.SetPtEtaPhiE(E_s/std::cosh(eta_s), eta_s, phi_s, E_s);
\end{lstlisting}

\paragraph{Six smearing variations (systematics).} The integer
\texttt{smear\_idx} ($0\text{--}5$) selects a systematic variation:
\texttt{base\_smear\_idx = idx\%3} picks one of three energy/position function
sets named in the code as
\texttt{global\_minimum}, \texttt{cE\_0p08}, \texttt{cE\_0p05}
(i.e.\ different assumed constant terms $c=0.08,\,0.05$); and for
\texttt{idx}$\ge 3$ an \emph{extra} 8\% constant term is added in quadrature:
\begin{lstlisting}
if (add_extra_const) E_res = std::sqrt(0.08f*0.08f + E_res*E_res);
\end{lstlisting}
These six histograms (\texttt{h\_m\_pt\_smear1..6}, filled at lines~740--745)
are how the smearing systematic is propagated.

\subsection{Where the smear functions come from}
The functions stored in \texttt{function\_compare\_wide.root} are produced by\\
\pp{/sphenix/u/bseidlitz/work/emcPerf/macros/figMaker/smear\_compare.C}.
This macro loads MC $\pi^0/\eta$ mass-vs-$p_T$ histograms
(\texttt{h\_m\_pt\_smear12/13/14}), fits each with the \texttt{MesonFitter}
class, and compares the fitted mean ($M_{\rm fit}/m_{\rm PDG}$, the
\emph{non-linearity}) and relative width ($\sigma/M_{\rm fit}$, the
\emph{resolution}) across smear variations labelled
\texttt{A(-1\%)/$\sqrt{E}$}, \texttt{A/$\sqrt{E}$}, \texttt{(A+1\%)/$\sqrt{E}$}.
The per-$p_T$ spread (max$-$min)/average is taken as the smearing systematic.
\textbf{Loop closed:} data widths $\Rightarrow$ choose smear function
$\Rightarrow$ store \texttt{TF1} $\Rightarrow$ apply in \texttt{smearPhot}
$\Rightarrow$ re-fit smeared MC $\Rightarrow$ verify it matches data.

\subsection{Related: MC energy-scale tuning}
Smearing fixes the \emph{width}; a separate small \emph{scale} correction fixes
the \emph{mean}. That lives in
\pp{/sphenix/u/bseidlitz/work/adjMCCalib/}: \texttt{convSwitch.C} calls
\texttt{scaleCDBTree(...)} to rescale the MC tower ADC$\to$E calibration
(\texttt{CEMC\_MC\_calib\_default\_v3.root}). Scale tuning and width smearing
together align MC to data in both peak position and peak width.

%==============================================================
\section{Unfolding --- what it is and how it is done here}
%==============================================================

\subsection{Teaching: the meaning of unfolding}
A measured spectrum is the \emph{true} spectrum convolved with detector effects
(resolution/smearing, efficiency, bin migration). \textbf{Unfolding} is the
statistical inversion that recovers the truth-level distribution from the
measured one. Formally, with a \emph{response matrix} $R_{ij}=P(\text{measured
bin }i\,|\,\text{truth bin }j)$ built from MC,
\[
  \vec{m}=R\,\vec{t}\quad\Longrightarrow\quad
  \vec{t}_{\rm unfold}=\mathcal{U}(R,\vec{m}),
\]
where naive matrix inversion amplifies statistical noise, so iterative
\textbf{Bayesian unfolding} (D'Agostini) is used instead --- regularised by a
small number of iterations. The MC that builds $R$ must already be
\emph{smeared} to match data resolution; otherwise the response is wrong and
the unfolded result is biased. This is the direct link from the smearing work
to PPG12.

\subsection{Toy unfolding and the resolution-mismatch bias}
\pp{/sphenix/u/bseidlitz/work/ppg12/toymcunfold/macro/resolutionsim.C}
is a pure teaching toy. It loads \texttt{libRooUnfold.so} (from
\texttt{/sphenix/user/egm2153/...}) and, for several pairs of
(MC resolution, data resolution), builds a response from a smeared exponential
$p_T$ spectrum and unfolds:
\begin{lstlisting}
RooUnfoldResponse *response = new RooUnfoldResponse(nbins,xlow,xup);
double recophotonpT = photonpT * r->Gaus(1, reso_mc);   // build response (MC reso)
response->Fill(recophotonpT, photonpT);
...
double recophotonpT = photonpT * r->Gaus(1, reso_data); // "data" with a DIFFERENT reso
RooUnfoldBayes *unfold = new RooUnfoldBayes(response, hMeasPT, 2); // 2 iterations
TH1D *hRecoPT = (TH1D*)unfold->Hreco();
\end{lstlisting}
It then plots the \textbf{bias} $(\text{unfolded}-\text{truth})/\text{truth}$.
The pairs \texttt{reso = \{(0.1,0.1),(0.1,0.11),(0.1,0.09),(0.05,0.06),
(0.05,0.04)\}} deliberately mismatch MC and data resolution to show how a
poorly-tuned smearing inflates the unfolding bias --- the quantitative argument
for why smearing must be tuned first.

\subsection{Realistic unfolding with Miss/Fake/Match}
\pp{/sphenix/u/bseidlitz/work/ppg12/toymcunfold/macro/jet\_pt\_unfolding.C}
is the full machinery on real JetValidation trees. It demonstrates the
bookkeeping that the toy omits:
\begin{itemize}[nosep]
  \item \texttt{response->Fill(reco,truth)} for matched ($\Delta R<0.3$) pairs;
  \item \texttt{response->Miss(truth)} for truth jets with no reco match;
  \item \texttt{response->Fake(reco)} for reco jets with no truth match;
  \item a half-sample \emph{closure test} (\texttt{resp\_half},
        \texttt{hMeasPTHalf}) to verify the procedure returns truth on
        independent MC;
  \item \texttt{RooUnfoldBayes(resp\_full,hMeasPT,1)} with 1--4 iterations as
        the regularisation parameter.
\end{itemize}
The PPG12 \texttt{README.md} confirms this is the production strategy:
``Bayesian unfolding (RooUnfold)'', with energy-resolution \emph{smearing
variations} (\texttt{config\_eres.yaml}) and unfolding iteration count listed
explicitly among the systematic uncertainties.

%==============================================================
\section{How the three concepts fit together}
%==============================================================
\begin{enumerate}[nosep]
  \item \textbf{Resolution} is measured in data as $\sigma(M)/M$ vs $p_T$ of the
        $\pi^0/\eta$ peaks (\texttt{resolutionStudy/}, \texttt{figMaker/}).
  \item \textbf{Smearing} injects exactly that much extra blur into MC so
        simulated peaks match data, via \texttt{pi0EtaByEta::smearPhot} using
        TF1s from \texttt{function\_compare\_wide.root}; six variations give the
        smearing systematic. Scale tuning (\texttt{adjMCCalib/}) aligns the mean.
  \item \textbf{Unfolding} uses that smeared MC to build the response matrix and
        Bayesian-inverts the measured spectrum back to truth for the final
        cross-section (\texttt{ppg12/toymcunfold/}, PPG12 pipeline).
\end{enumerate}

\noindent\textbf{One-line glossary.}
\emph{Resolution} = how blurry the detector is (measured).
\emph{Smearing} = adding that blur to MC (applied).
\emph{Unfolding} = removing the blur from a measured spectrum to recover truth
(corrected).

\bigskip
\noindent\footnotesize All inspected files are read-only collaborator areas
(owner \texttt{bseidlitz}). Exact paths, file sizes, and line references are
recorded in \texttt{RESO\_SMEAR\_UNFOLD\_LOG.txt} in this directory.

\end{document}
